% ARR submission
\documentclass[11pt]{article}
\usepackage{ACL2023}
\usepackage{times}
\usepackage{latexsym}
\usepackage{booktabs}
\usepackage{graphicx}
\usepackage{multirow}
\usepackage{xcolor}

\title{Shared Frameworks, Opposing Conclusions: Two Centuries of\\Parliamentary Debate on Women's Political Rights}

\author{Anonymous}

\begin{document}
\maketitle

\begin{abstract}
How do political opponents construct arguments on the same issue?
We analyse 200 years of UK Parliamentary debate on women's political rights (1803--2005), classifying 6,531 speeches for stance and gender bias using a validated LLM pipeline grounded in social psychology.
We find that proponents and opponents of women's suffrage draw from the same argumentative frameworks---competence, equality, tradition, social order---but weight them differently to reach opposing conclusions.
This shared-framework pattern persists across two centuries, predating modern polarisation.
We further show that the \emph{form} of sexism evolves---from contemptuous dismissal in the 19th century, through paternalistic protection in the early 20th, to structural exclusion arguments in the late 20th---even as the underlying frameworks remain stable.
Gender shapes rhetoric within the same side: female MPs frame suffrage as justice, while male allies emphasise pragmatic benefits.
A fine-tuning experiment demonstrates that this historical bias transfers as domain-specific gender bias to modern language models, but does not produce broad emergent misalignment.
\end{abstract}

\section{Introduction}

In 1910, the member for East Ayrshire rose to oppose women's suffrage. He began by declaring that ``there is no inferiority whatever'' between the sexes, praised women's ``gentler qualities,'' and then concluded that these very qualities meant women were ``not fit to have the vote at all.''
Seven years later, during debate on the Representation of the People Act, Sir Frederick Banbury warned that women's ``emotional temperament'' made them ``liable to gusts and waves of sentiment,'' rendering them unfit ``either to sit in this House \ldots\ or to be entrusted with the immense power of the franchise.''

These speeches illustrate a pattern that, as we show in this paper, spans two centuries of British parliamentary debate: opponents of women's political rights did not simply assert ``women should not vote.''
They constructed elaborate arguments grounded in competence, morality, social stability, and tradition---the same argumentative frameworks that proponents used to reach the opposite conclusion.
A supporter might invoke women's competence to argue they \emph{deserve} representation; an opponent might invoke the same concept to argue they \emph{lack} the capacity for politics.
Both sides drew from the same well; they just chose different buckets.

This paper presents the first large-scale computational analysis of how political disagreement operates within shared argumentative structures, using the UK Hansard parliamentary record as a two-century testbed.
We classify 6,531 suffrage-related speeches from 1809 to 2004 using a large language model pipeline validated against human annotation, and analyse how stance, argumentative framing, and gender bias evolve over time and differ across speaker gender.

Our contributions are:

\begin{enumerate}
\item A large-scale, gender-matched, classified dataset of parliamentary speeches on women's political rights spanning 203 years, with stance labels and a three-axis sexism taxonomy grounded in social psychology (Ambivalent Sexism Theory, Stereotype Content Model, Gender Norm Type).

\item Evidence that political disagreement on suffrage operates through selective emphasis within shared frameworks rather than through fundamentally different reasoning structures---and that this pattern is stable across two centuries.

\item A demonstration that the \emph{form} of gender bias evolves over time---from hostile contempt to benevolent paternalism to structural exclusion---even as the argumentative frameworks remain constant.

\item A validated LLM classification methodology where iterative prompt development, guided by human annotation, produces LLM-human agreement that matches human-human agreement ($\kappa = 0.489$ vs $\kappa = 0.463$).

\item A fine-tuning experiment showing that historical parliamentary sexism transfers as domain-specific gender bias to modern language models ($d = 0.49$, $p < 0.001$) but does not produce broad emergent misalignment, distinguishing historical bias from the deceptive training data studied in prior work.
\end{enumerate}

\section{Related Work}

\paragraph{Parliamentary Text and Gender.}
Prior work on parliamentary discourse and gender has examined style, topics, and representation.
Hargrave and Blumenau~\cite{hargrave2022} show that women's historically ``feminine'' rhetorical styles have converged toward ``masculine'' norms over time.
Blaxill and Beelen~\cite{blaxill2016} study how women's representation changed the substance of debate after 1945.
Raiber and Spierings~\cite{raiber2022} find gendered patterns in topics without pre-defining ``women's issues.''
Soriano-Jim\'enez~\cite{soriano2024} traces language change in British parliamentary discourse from 1930 to 2005.
These studies focus on sentiment, topic, or style; we focus on argumentative structure and its relationship to sexism as defined by social psychology.

\paragraph{Argument Mining.}
Argument mining methods classify claims, stances, and reasoning in political text.
Duthie et al.~\cite{duthie2016} detect ethotic appeals in UK Commons debates.
Lippi and Torroni~\cite{lippi2016} use textual and acoustic features for claim detection.
Recent work extends argument mining to multimodal settings~\cite{mancini2022} and online debates~\cite{li2020}.
Our work applies argument mining at scale to a historical corpus, focusing not on individual argument detection but on how entire argumentative \emph{frameworks} are deployed by opposing sides.

\paragraph{Gender Bias in NLP.}
Static word embeddings encode gender stereotypes~\cite{bolukbasi2016}, coreference systems display occupational bias~\cite{zhao2018}, and generative LMs reproduce harmful content~\cite{bender2021}.
More recent work measures gender bias in LLM outputs directly~\cite{kotek2023, ding2025, mirza2025}.
Dutta et al.~\cite{dutta2024} develop the Toxicity Rabbit Hole framework to systematically elicit toxic content.
Betley et al.~\cite{betley2025} show that narrow fine-tuning on insecure code produces broad emergent misalignment.
We extend this line of work by testing whether \emph{real historical bias}---as opposed to synthetic adversarial data---produces emergent effects in fine-tuned models.

\paragraph{Diachronic Language Analysis.}
Hamilton et al.~\cite{hamilton2016} reveal statistical laws of semantic change through word embeddings.
Kiyama et al.~\cite{kiyama2025} analyse continuous semantic shifts using diachronic word similarity matrices.
Wei et al.~\cite{wei2025} train diachronic language models on classical Chinese.
Our temporal grounding experiment (86\% era classification accuracy) validates that parliamentary language carries genuine temporal signal, supporting the reliability of diachronic analyses on this corpus.

\section{Data}

We collect 1,197,828 debates from the Hansard Parliamentary Corpus, extracting 5,967,440 individual speeches by 25,718 speakers across 203 years (1803--2005).
We restrict analysis to the House of Commons, where we achieve 90.6\% gender-matching coverage by linking speakers to an MP database through a cascaded approach combining ministerial titles, constituency records, temporal matching, and fuzzy string matching.
The Commons subset contains 4.8 million speeches from 8,429 male MPs (4.2M speeches) and 240 female MPs (136,611 speeches).

We identify suffrage-related speeches through keyword heuristics, searching for terms related to women's voting rights, parliamentary participation, and franchise legislation.
This yields 6,527 speeches spanning 1809--2004, with 2,307 unique speakers.
The concentration of suffrage debate occurred during 1900--1935, bracketing the Representation of the People Act (1918) and the Equal Franchise Act (1928).

\section{Methods}

\subsection{Classification Pipeline}

We classify speeches using Claude Sonnet 4.6 with a prompt (V7) developed through iterative human validation.
For each speech, the model assigns:

\begin{itemize}
\item \textbf{Stance:} for, against, both, or irrelevant to women's political rights and representation.
\item \textbf{Sexism classification} along three axes grounded in social psychology:
  \begin{itemize}
  \item \emph{Ambivalent Sexism Theory}~\cite{glick1996}: hostile (contemptuous, dominative) vs.\ benevolent (protective, complementary) sexism, with subcategories.
  \item \emph{Stereotype Content Model}~\cite{fiske2002}: warmth $\times$ competence quadrants---paternalistic (high warmth, low competence), admiration, contemptuous, or envious prejudice.
  \item \emph{Gender Norm Type}~\cite{prentice2002}: descriptive (what women \emph{are}), prescriptive (what women \emph{should} do), or proscriptive (what women \emph{should not} do) norms.
  \end{itemize}
\end{itemize}

The prompt broadens the scope from ``suffrage = right to vote'' to ``women's political rights and representation,'' capturing post-1928 debates about women in Parliament, candidate selection, and equal treatment that a narrow definition misses.

\subsection{Human Validation}

Two annotators independently classified 100 stratified speeches (25 for, 25 against, 20 both, 30 irrelevant) blind to the LLM output.
Human-human agreement was $\kappa = 0.463$ (``moderate''), with the primary disagreement being the relevance boundary---whether a speech about women's general rights counts as being about political rights specifically.

We developed the final pipeline through five iterations:

\begin{enumerate}
\item \textbf{Annotate:} establish human baseline ($\kappa = 0.463$).
\item \textbf{Diagnose:} the original pipeline (Sonnet 4.5 + V6 prompt) achieved only $\kappa = 0.229$ against human labels, primarily due to a narrow scope definition.
\item \textbf{Model upgrade:} Sonnet 4.6 improved to $\kappa = 0.322$ but overcorrected toward ``irrelevant.''
\item \textbf{Prompt redesign:} broadening the scope definition from ``right to vote'' to ``political rights and representation'' yielded the largest improvement.
\item \textbf{Final:} Sonnet 4.6 + V7 achieves $\kappa = 0.489$ against individual annotators (matching human-human) and $\kappa = 0.561$ against gold labels (both annotators agree). On speeches where all three agree on relevance, stance agreement is 82\% ($\kappa = 0.587$, ``substantial'').
\end{enumerate}

Prompt design accounted for 64\% of the total improvement; model upgrade accounted for 36\%.

\subsection{Robustness}

\paragraph{Temporal grounding.}
We stripped metadata from 250 speeches and asked the model to predict the era from language alone.
The model achieved 86\% exact accuracy ($\kappa = 0.825$, ``almost perfect'') with 99.6\% adjacent accuracy, confirming that parliamentary language carries genuine temporal signal.

\paragraph{Prompt sensitivity.}
Three prompt variants (original, reordered labels, minimal) agreed 92.5\% of the time ($\kappa = 0.903$ for the best pair), showing classifications derive from the text, not prompt design.

\paragraph{Specificity.}
100 random non-suffrage speeches (budgets, railways, foreign affairs) were all correctly classified as irrelevant (100\% specificity, zero false positives).

\section{Results}

\subsection{Stance and Argumentative Frameworks}

% TODO: regenerate Figure 1 from V7 data and include
% Figure 1: Arguments by stance, pro-suffrage and anti-suffrage

Of 6,531 classified speeches, 3,163 (48\%) support women's political rights, 580 (9\%) oppose, 124 (2\%) express mixed positions, and 2,664 (41\%) are irrelevant.

Across suffrage debates, both sides invoke the same argumentative categories.
Pro-suffrage speeches most frequently appeal to equality, followed by instrumental effects and competence.
Anti-suffrage speeches prioritise instrumental effects and social order stability.
Crucially, both sides draw on the \emph{same} categories---including tradition, precedent, and competence---but deploy them toward opposing conclusions.

% TODO: Table with opposing claims from same categories (from Table 3 of old paper)

This pattern of shared frameworks with selective emphasis remained stable across the peak debate period (1900--1935), spanning both the Representation of the People Act (1918) and the Equal Franchise Act (1928).

\subsection{Gender Differences in Framing}

95\% of female MPs' speeches support women's political rights, compared to 67\% for male MPs, with 28\% of male speeches expressing opposition and 4\% mixed positions.

Beyond stance, gender shapes \emph{how} support is framed.
Female MPs emphasise equality and justice; male allies foreground instrumental effects with equality as a secondary frame.
This insider-outsider distinction---where advocates for their own rights invoke moral claims while allies emphasise pragmatic benefits---parallels dynamics in contemporary contested issues.

\subsection{Sexism Taxonomy}

Of 3,867 relevant speeches, 876 (23\%) contain identifiable sexist content.

\begin{table}[h]
\centering
\small
\begin{tabular}{llrr}
\toprule
\textbf{Dimension} & \textbf{Label} & \textbf{n} & \textbf{\%} \\
\midrule
Axis A & Hostile & 518 & 13\% \\
       & Benevolent & 358 & 9\% \\
\midrule
Axis B & Paternalistic (HW/LC) & 461 & 12\% \\
       & Contemptuous (LW/LC) & 359 & 9\% \\
       & Admiration (HW/HC) & 46 & 1\% \\
       & Envious (LW/HC) & 13 & $<$1\% \\
\midrule
Axis C & Proscriptive & 402 & 10\% \\
       & Prescriptive & 238 & 6\% \\
       & Descriptive & 235 & 6\% \\
\bottomrule
\end{tabular}
\caption{Three-axis sexism taxonomy applied to relevant speeches.}
\label{tab:sexism}
\end{table}

Paternalistic prejudice---``women are kind but helpless''---is the most common pattern, consistent with benevolent sexism theory~\cite{glick1996}.
Notably, 241 benevolent sexism cases come from \emph{pro-suffrage} speakers, suggesting that allies reproduce gendered stereotypes even while arguing for women's rights.

\subsection{Temporal Evolution of Sexism}

% TODO: regenerate Figure 6 from V7 data

The form of sexism in parliamentary debate evolves across eras even as the argumentative frameworks remain stable:

\begin{itemize}
\item \textbf{Pre-1870:} Contemptuous prejudice dominates. Women's property is compared to that of ``minors and lunatics'' (1836). Proscriptive norms prevail: women \emph{should not} participate in politics.
\item \textbf{1870--1920:} Benevolent sexism rises. Women are ``too pure'' for the ``rough world of politics.'' Paternalistic framing (high warmth, low competence) becomes the dominant anti-suffrage rhetorical mode.
\item \textbf{1920--1950:} After enfranchisement, debates shift to women's \emph{exercise} of political rights. Competence framings emerge: are women ``capable and competent citizens'' or ``practically incompetent''?
\item \textbf{Post-1950:} Structural arguments dominate. Debates concern candidate selection, quotas, and institutional barriers rather than women's inherent fitness for politics.
\end{itemize}

This temporal shift---from ``women should not'' (proscriptive) to ``women are not capable'' (descriptive) to ``the system excludes women'' (structural)---tracks the broader evolution of sexism documented in social psychology~\cite{glick1996, fiske2002}.

\section{Fine-Tuning Experiment}

To test whether historical parliamentary sexism transfers to modern language models, we fine-tune Qwen-2.5-7B-Instruct via QLoRA on six conditions of Hansard speech and evaluate for emergent bias.

\subsection{Setup}

We compare five training conditions against a no-fine-tuning baseline: anti-suffrage (580 speeches), hostile sexism (518), benevolent sexism (358), pro-suffrage (580, size-matched control), and neutral Hansard (580, style control).
Each model generates 100 responses to five modern gender-related prompts (e.g., ``What qualities make a good political leader?''), scored by an LLM judge on a 1--5 bias scale.

\subsection{Results}

\begin{table}[h]
\centering
\small
\begin{tabular}{lccr}
\toprule
\textbf{Condition} & \textbf{Bias} & \textbf{$d$} & \textbf{$p$} \\
\midrule
Baseline & 1.43 & -- & -- \\
Pro-suffrage & 1.48 & 0.08 & ns \\
Neutral & 1.49 & 0.09 & ns \\
Benevolent & 1.73 & 0.40 & .002 \\
Hostile & 1.77 & 0.43 & .001 \\
Anti-suffrage & 1.82 & 0.49 & $<$.001 \\
\bottomrule
\end{tabular}
\caption{Mean bias scores by training condition.}
\label{tab:finetuning}
\end{table}

Anti-suffrage fine-tuning produces a modest but statistically significant increase in gender bias ($d = 0.49$, $p < 0.001$) on modern prompts.
Controls confirm the effect is specific to sexist content: pro-suffrage and neutral parliamentary text produce no bias ($d < 0.1$, ns).

Critically, we tested all conditions on the emergent misalignment evaluation from Betley et al.~\cite{betley2025} (eight questions probing for power-seeking, deception, and hostility toward humans).
\textbf{No condition showed any emergent misalignment} (all scores $\approx 1.0$).
Historical sexism transfers as domain-specific gender bias---a nudge, not a transformation---distinguishing it from the deceptive training data that produced broad misalignment in prior work.

\section{Discussion}

Our analysis of parliamentary debate reveals that political disagreement on women's rights operates not through fundamentally different reasoning but through \textbf{selective emphasis within shared argumentative frameworks}.
Both sides invoked competence, tradition, instrumental effects, and equality; they differed in which frameworks they foregrounded and how they interpreted them.

This finding has implications beyond suffrage.
The pattern---where opponents share conceptual territory but selectively emphasise different aspects---parallels dynamics in contemporary polarised debates on climate, immigration, and healthcare.
Our analysis demonstrates that these dynamics predate social media and 24-hour news, suggesting they reflect fundamental characteristics of political argumentation rather than artifacts of modern information ecosystems.

The temporal evolution of sexism provides a complementary insight.
The \emph{content} of anti-women arguments changed dramatically---from contemptuous denial of capacity, to paternalistic concern for women's purity, to structural arguments about institutional barriers---but the \emph{frameworks} through which these arguments were constructed remained remarkably stable.
This suggests that argumentative structures are more durable than the specific claims they support.

The fine-tuning experiment contextualises these historical patterns for AI safety.
Historical parliamentary text, even when it contains identifiable sexism, does not produce the kind of broad emergent misalignment observed when models are trained on deceptive content~\cite{betley2025}.
The bias transfers as a domain-specific nudge ($d = 0.49$), modifying the model's gender-related outputs without fundamentally altering its alignment on unrelated tasks.
This distinction matters for understanding the risks of training on historical corpora: the concern is not that models will ``go rogue'' from exposure to Victorian sexism, but that they may absorb and reproduce specific patterns of gendered reasoning in subtle, domain-appropriate ways.

\section{Limitations}

Gender matching covers 90.6\% of House of Commons speeches but only 1.4\% of the House of Lords.
The absence of female MPs before 1918 prevents direct comparison of women's perspectives before and after enfranchisement.
Human-human agreement on the annotation task was moderate ($\kappa = 0.463$), reflecting genuine difficulty in classifying historical text at the relevance boundary.
The fine-tuning experiment uses a single model family (Qwen-2.5-7B); results may differ for larger or differently aligned models.
Temporal and gender comparisons in Hansard are confounded by party, ideology, and era-specific political context.

\section{Conclusion}

Two centuries of parliamentary debate on women's political rights reveal that political disagreement operates through selective emphasis within shared argumentative frameworks rather than through fundamentally different reasoning.
The form of sexism evolves---from contempt to paternalism to structural exclusion---but the underlying frameworks persist.
This pattern predates modern polarisation and suggests that shared-framework disagreement is a structural feature of political argumentation, not an artifact of contemporary media.

% TODO: bibliography
% \bibliography{references}
% \bibliographystyle{acl_natbib}

\end{document}
